% Generated draft: requires mathematical and source review.
\documentclass[11pt]{article}
\usepackage[margin=25mm]{geometry}
\usepackage{fontspec}
\setmainfont{Latin Modern Roman}
\usepackage{amsmath,amssymb,mathtools,unicode-math}
\setmathfont{Latin Modern Math}
\usepackage{enumitem,needspace,xurl,hyperref}
\setlength{\parindent}{0pt}
\setlength{\parskip}{5pt}
\setlength{\emergencystretch}{3em}
\newcommand{\N}{\mathbb{N}}
\newcommand{\Z}{\mathbb{Z}}
\newcommand{\Q}{\mathbb{Q}}
\newcommand{\R}{\mathbb{R}}
\newcommand{\C}{\mathbb{C}}
\newcommand{\F}{\mathbb{F}}
\newcommand{\A}{\mathbb{A}}
\newcommand{\PP}{\mathbb{P}}
\newcommand{\E}{\mathbb{E}}
\newcommand{\eps}{\varepsilon}
\newcommand{\dd}{\,\mathrm{d}}
\begin{document}
\section*{UnsolvedMath ID: 30001005}
\textbf{Source problem code:} catalogue entry
\subsection{Definitions and mathematical statement}
Let \(n\geq7\), and let \(M_0\subset\R^{n+1}\) be a hypersurface whose singularities are isolated points modeled on regular hypercones. A hypercone here means an \(n\)-dimensional cone invariant under positive dilations; the usual regularity convention is smoothness away from its vertex. The extraction does not state minimality, a precise asymptotic regularity condition, or a convergence topology for the conical models.

Let \((M_t)_{t\geq0}\) be a mean-curvature flow with initial data \(M_0\), in the intended weak-flow sense at the singular time. Where smooth, mean-curvature flow has normal velocity equal to its mean-curvature vector. The question asks both:
\begin{enumerate}
\item Near each singular point \(p\), does the flow have asymptotically self-similarly expanding behavior as \(t\downarrow0\)? The expected parabolic scale is \(t^{-1/2}(M_t-p)\); convergence to an expanding profile would need its topology, multiplicity, and selection conditions specified. A self-similar expanding profile produces the family \(p+\sqrt t\,\Sigma\) for \(t>0\) under the chosen time normalization.
\item Does the flow become smooth for a short positive time, that is, is \(M_t\) smooth for \(0<t<\varepsilon\) for some \(\varepsilon>0\) in the intended local or global sense?
\end{enumerate}
The cleaned question restricts the remaining target to \(n\geq7\). The original asks without this restriction and says that affirmative results hold for \(n\leq6\). Which weak flow is meant, whether assertions concern every such flow or an appropriate selected flow, and whether smoothing is global or local near each singularity are not specified.
\subsection{Short English statement}
In hypersurface dimensions \(n\geq7\), ask whether mean-curvature flow from isolated regular conical singularities immediately smooths and approaches expanding self-similar profiles near those points. The weak-flow choice and precise asymptotic and smoothing statements still require hypotheses.
\subsection{Sources}
\begin{itemize}
\item Local UnsolvedMath registry, permanent ID \(30001005\), source code \texttt{catalogue entry}, ``Mean-Curvature Flow from Isolated Conical Singularities.'' The supplied assessment attributes the questions to Ilmanen. Registry curation: Ulam AI / UnsolvedMath Contributors, CC BY 4.0; underlying sources retain their supplied provenance.
\item \emph{Geometrie} (2009), Oberwolfach Reports, pp.\ 1933--1978, DOI \url{https://doi.org/10.4171/owr/2008/34}.
\item Supplied ``mean-curvature flow from conical singularities'': \url{https://doi.org/10.1007/s00222-024-01296-8}.
\item Supplied literature triage, \texttt{literature\_checked\_at}: 2026-08-21: labeled partially solved; says Ilmanen's questions are affirmative for \(n\leq6\), and Chodosh--Daniels-Holgate--Schulze proved short-time smoothness for isolated conical singularities in that range, while the unrestricted dimensional assertion remains open. This is supplied dated curation, not independently verified; the full hypotheses of the cited results require review.
\end{itemize}
\end{document}
